%% ma: L12-B14-0118
%% lop: 12
%% chuong: 5
%% bai: 14
%% dang: 12.14.28
%% loai: DS
%% muc: [NB, TH, TH, VD]
%% dap_an: "ĐSSS"
%% nguon: "Ảnh giáo viên gửi 07/11/2026 (Tư Duy Mở - tuduymo.com, đề tự luyện Phương trình mặt phẳng, đề khoảng cách), Câu 14"
%% kien_thuc: "Bài 14 (hai mặt phẳng song song, khoảng cách từ một điểm đến mặt phẳng), phương trình mặt phẳng tổng quát"
%% trang_thai: chinh-thuc
%% ghi_chu: "Tài liệu gốc không kèm đáp án; đáp án do Claude giải và kiểm bằng sympy. Dạng đã duyệt 07/11/2026."
\begin{ex}
Trong không gian $Oxyz$, cho hai điểm $A(-3;0;3)$, $B(a;b-1;2a+2b-3)$, với $a$ và $b$ là các số thực thay đổi.
\choiceTF
{\True Phương trình tổng quát của mặt phẳng đi qua $A$ và song song với $(Oxy)$ là $z-3=0$.}
{Điểm $B$ thuộc một mặt phẳng $(\alpha)$. Phương trình tổng quát của $(\alpha)$ là $2x+2y-z=0$.}
{$d\left(A,(\alpha)\right)=3$.}
{Phương trình tổng quát của mặt phẳng $(\beta)\parallel(\alpha)$ sao cho $d\left((\alpha),(\beta)\right)=2$ tương ứng là $(\beta)\colon 2x+2y-z+12=0$.}
\loigiai{a) Mặt phẳng song song $(Oxy)$ có dạng $z=c$; qua $A$ nên $z-3=0$. Đúng.

b) $x=a$, $y=b-1$, $z=2x+2(y+1)-3=2x+2y-1$, nên $(\alpha)\colon 2x+2y-z-1=0$. Sai.

c) $d\left(A,(\alpha)\right)=\dfrac{|-6-3-1|}{3}=\dfrac{10}3$. Sai.

d) $(\beta)\colon 2x+2y-z+D=0$, $\dfrac{|D+1|}{3}=2\Rightarrow D=5$ hoặc $D=-7$. Sai.}
\end{ex}
